\documentclass[11pt]{article}
\usepackage{amsmath,amssymb,amsthm,hyperref}
\newtheorem{theorem}{Theorem}
\newtheorem{corollary}[theorem]{Corollary}
\title{Residual coordinates and shared-deformation barriers for circuit counting}
\author{Astra P3 research construction}
\date{27 October 2026}
\begin{document}
\maketitle
\paragraph{Scope.} Exact restrictions on proposed extensions of the supplied
counting construction; not lower bounds for arbitrary circuit algorithms or
claims of historically new algebra.

\section{Gate residuals are independent coordinates}
Let an acyclic Boolean circuit have independent inputs $x\in\mathbb F_2^n$
and $s$ gates in topological order. For stored gate values $y_g$, define
\[
 e_g=y_g+f_g(x,y_1,\ldots,y_{g-1})\quad\text{in }\mathbb F_2.
\]
The Boolean operation $f_g$ may have arbitrary arity.
\begin{theorem}[Residual-coordinate bijection]
The map $(x,y)\mapsto(x,e)$ is a bijection. At every fixed input, every
residual vector $e\in\mathbb F_2^s$ corresponds to exactly one trace.
\end{theorem}
\begin{proof}
Its inverse successively computes $y_g=e_g+f_g(x,y_1,\ldots,y_{g-1})$.
\end{proof}
\begin{corollary}[Parity consistency filters]
For a binary matrix $A$ of rank $r$, the condition $Ae=0$ admits exactly
$2^{n+s-r}$ traces. It selects only the correct trace at every input iff $r=s$.
\end{corollary}
\begin{proof}
The kernel has $2^{s-r}$ elements; the zero residual is the correct trace.
\end{proof}
Clause-violation vectors in Section 6 of the governing paper have at most
$2^n$ possible values on the original input cube. Freely summed gate residuals
take all $2^s$ values even at a fixed input. The short-test argument cannot
silently equate these domains. Eliminating gate variables first remains possible.

\paragraph{Normalization does not preserve arbitrary observables.}
For two constant-zero gates, the filter $e_1=e_2$ admits $00,11$.
The observable $y_2$ has correct value zero but filtered sum one. No nonzero
global normalization corrects this and the constant-one observable together.
This elementary example is not a hardness statement.

\section{Nonlinear residual filters have a degree budget}
\begin{theorem}[Exact zero-selection degree]
If $q_1,\ldots,q_a\in\mathbb F_2[e_1,\ldots,e_s]$ have common Boolean
zero set exactly $\{0\}$, then $\sum_j\deg q_j\ge s$.
\end{theorem}
\begin{proof}
On the Boolean cube, $\prod_j(1+q_j(e))$ is the zero-set indicator.
Its unique multilinear reduction is $\prod_i(1+e_i)$, of degree $s$.
Multilinear reduction never increases degree, while the original product has
degree at most $\sum_j\deg q_j$.
\end{proof}
This concerns tests in residual coordinates. It does not transfer unchanged
through a nonlinear substitution or an output-specific simplification.

\begin{theorem}[Signed universal selector]
If a rational polynomial $w$ satisfies
$\sum_{e\in\{0,1\}^s}w(e)G(e)=G(0)$ for every Boolean observable $G$,
then $w(0)=1$, $w(e)=0$ for $e\ne0$, and $\deg w\ge s$.
\end{theorem}
\begin{proof}
Choose each point indicator for $G$. The multilinear representative of $w$
is then $\prod_i(1-e_i)$, of degree $s$.
\end{proof}
Signed cancellation may help a selected output without being a universal
low-degree consistency filter. Such a method must construct and pay for its
output-specific information.

\section{A shared linear bottleneck has a rank-loss budget}
\begin{theorem}
Let $P(T)$ be an $n\times n$ polynomial matrix over a characteristic-zero
field, with entry degrees at most $R$ and $P(0)=I$. At distinct nonzero
points $t_1,\ldots,t_Q$, put $b_j=n-\operatorname{rank}P(t_j)$. Then
\[
 \sum_j b_j\le nR.
\]
If every point has rank at most $s$, then $R/Q\ge1-s/n$.
\end{theorem}
\begin{proof}
The nonzero determinant has degree at most $nR$. Constant invertible row
operations make $b_j$ rows vanish at $t_j$. Their polynomial entries are
divisible by $T-t_j$, so the determinant is divisible by $(T-t_j)^{b_j}$.
Multiply these divisibilities at distinct points and compare degrees.
\end{proof}
This rules out restoring all inputs at zero while passing them through a
small common linear image at many cheap nodes with $R/Q=o(1)$. It does not
rule out full-rank block-local maps or the paper's incidence-specific forms.

\section{Executed checks}
The exact-integer program \texttt{residual\_fibers.py} checks 32 acyclic
circuits, $1\le n\le4$, $1\le s\le8$. All 15,300 input/residual pairs pass
the inverse-map identity. Random-rank filters admit $2^{n+s-r}$ traces.
Boolean M\"obius inversion gives zero-selector degree $s$ in every case.
The two-gate observable counterexample passes. These finite checks supplement
the proofs and are not a practical performance benchmark.

\paragraph{Source.} OpenAI, \emph{The Exponential Time Hypothesis Is False},
revision 25 October 2026, Sections 3--6 and Appendix D, commit
\texttt{7ef5c0cc5ca4b0e072b7a050fca09ab24bc6fa3a}.
The supplied PDF matches Git blob
\texttt{0503566e95bc587a4076573a45132d97531b8ba6} byte-for-byte.
\end{document}
